\documentclass[sn-mathphys, oneside]{sn-jnl}
\usepackage[T5,T1]{fontenc}
\usepackage[utf8]{inputenc}
\usepackage[english]{babel}
\newcommand{\vn}[1]{{\fontencoding{T5}\selectfont #1}}
\newcommand{\vntt}[1]{{\fontencoding{T5}\selectfont\texttt{#1}}}
\usepackage{indentfirst}
% Used to describe data
\usepackage{listings}
\usepackage{xcolor}
\lstdefinelanguage{json}{
    basicstyle=\ttfamily\footnotesize,
    numbers=left,
    numberstyle=\tiny,
    stepnumber=1,
    breaklines=true,
    frame=single,
    literate=
     *{0}{{{\color{blue}0}}}{1}
      {1}{{{\color{blue}1}}}{1}
      {2}{{{\color{blue}2}}}{1}
      {3}{{{\color{blue}3}}}{1} 
}      
%
\usepackage{hyperref}
\usepackage{url}
\usepackage{hyperref}
\usepackage{placeins}
\usepackage{graphicx}
\usepackage{amsmath}
\usepackage{booktabs}
\usepackage[font=small,labelfont=bf,labelsep=colon]{caption}
\captionsetup[table]{textfont=it, name=Table, labelsep=colon, skip=6pt}
\usepackage{booktabs}
\usepackage{algorithm}
\usepackage{algpseudocode} 
\usepackage{array}
\usepackage{tabularx}
\usepackage[numbers]{natbib}
\usepackage{amssymb}
\usepackage{pifont}
\usepackage{float}
\usepackage{mathptmx} 
\usepackage[most]{tcolorbox}
% \documentclass{article}
% \usepackage[utf8]{inputenc}
% \usepackage[vietnamese]{babel}
\usepackage{tikz}
\usepackage{subcaption}
\usepackage{graphicx}
\usepackage{tocloft} % Package for customizing the table of contents

\usepackage{algorithm}
\usepackage{algpseudocode}

\makeatletter
\newenvironment{breakablealgorithm}
  {% \begin{breakablealgorithm}
   \begin{center}
     \refstepcounter{algorithm}%
     \hrule height.8pt depth0pt \kern2pt%
     \renewcommand{\caption}[2][\relax]{%
       {\raggedright\textbf{\ALG@name~\thealgorithm} ##2\par}%
       \ifx\relax##1\relax
         \addcontentsline{loa}{algorithm}{\protect\numberline{\thealgorithm}##2}%
       \else
         \addcontentsline{loa}{algorithm}{\protect\numberline{\thealgorithm}##1}%
       \fi
       \kern2pt\hrule\kern2pt
     }
  }{%
     \kern2pt\hrule\relax%
   \end{center}
  }
\makeatother

% 1. Prepend "Figure" before the number
\renewcommand{\cftfigpresnum}{Figure } 

% 2. Increase the space reserved for the figure number
% (Without this, "Figure 1" label would overlap with the caption)
\setlength{\cftfignumwidth}{4.5em} 


\usetikzlibrary{shapes, arrows.meta, positioning, calc, fit, shadows}
\newcommand{\pending}[1]{\textbf{\textcolor{red!65!black}{[PLACEHOLDER: #1]}}}
\newcommand{\artifact}[1]{\nolinkurl{#1}}
% 
% --- INSERT THIS BLOCK INTO THE PREAMBLE (AFTER PACKAGE IMPORTS) ---
\geometry{
    paper=a4paper,      % A4 paper size
    left=2.5cm,         % Left margin 2.5cm
    right=2.5cm,        % Right margin 2.5cm (matches left margin)
    top=2.5cm,          % Top margin
    bottom=2.5cm,       % Bottom margin
    bindingoffset=0cm   % Important: keep binding offset at 0 to avoid text shift
}
\raggedbottom
% -----------------------------------------------------------
% Start a new page for the main content
\begin{document}

% --- Title page ---
\begin{titlepage}
    \centering
    
    % --- Header Section: Logo and Ministry of Education and Training ---
    \begin{minipage}{0.45\textwidth}
        \begin{flushleft}
            % Make sure logo.png has proper proportions, or adjust width accordingly
            \includegraphics[width=5cm]{logo.png} 
        \end{flushleft}
    \end{minipage}
    \hfill
    \begin{minipage}{0.5\textwidth}
        \begin{flushright}
            \textbf{MINISTRY OF EDUCATION AND TRAINING}
        \end{flushright}
    \end{minipage}

    \vspace{2.5cm}

    % --- University name ---
    {\Huge \textbf{FPT UNIVERSITY}\par}

    \vspace{2.5cm}

    % --- Document type ---
    {\huge Capstone Project Document \par}

    \vspace{0.5cm}
    % % Horizontal rule
    % \rule{\textwidth}{0.4pt}
    % \vspace{0.5cm}

    % --- Thesis title ---
    {\Large \textbf{VaxiS: An Evidence-Grounded Retrieval-Augmented Generation System for Vietnamese Vaccine Information Retrieval} \par}

    \vspace{2.5cm}

    % --- 4. INFORMATION TABLE (framed as in the reference image) ---
    \begin{center}
        \renewcommand{\arraystretch}{1.5} % Increase row height for readability
        \setlength{\arrayrulewidth}{0.8pt} % Table border thickness
        
        % Create table: Column 1 (4.5cm), Column 2 (10cm)
        \begin{tabular}{|p{4.5cm}|p{10cm}|}
            \hline
            % Group heading row (merged cells)
            \multicolumn{2}{|c|}{\textbf{AIP491 GROUP 5}} \\ 
            \hline
            \textbf{Group Members} & 
           Nguyen Manh Hung - HE194092  \newline
           Tran Trong Nghia - HE190336
            \newline
            Ho Minh Hieu - HE194091 \\
            \hline
            \textbf{Supervisor} &  Mr Tran Van Ha\\ % Mr
            \hline
            \textbf{Capstone Project code} & AIP491  \\
            \hline
        \end{tabular}
    \end{center}

    \vfill

    % --- Footer: Location and date ---
    {\large - Hanoi, 11/2026 - \par}
    
\end{titlepage}

\clearpage % new page for the article body following sn-jnl
%=====================================
% 2) Front matter following sn-jnl convention
%=====================================
\begin{center}
    \bfseries\Large Abstract
\end{center}


Vietnamese vaccine information is published in documents that vary in structure and are often written for specialist readers. VaxiS is a Vietnamese-language retrieval-augmented generation system designed to help users locate relevant passages and trace an answer to its documentary evidence. The repository contains a structure-aware Markdown indexing pipeline, lexical and dense retrieval paths, and a hybrid retrieval pipeline that combines dense and learned sparse rankings before optional reranking. A shared generator uses the retrieved chunks as context and tracks citations to chunk identifiers. The experimental design follows six configurations in order: Naive, Advanced, Hybrid (baseline), Table-aware, Graph RAG, and Agent RAG. The working benchmark contains 300 Vietnamese questions spanning single-hop, multi-hop, and table-reasoning tasks. \pending{For the final three configurations, implementation and evaluation documentation are pending. Insert the verified comparison findings, with the benchmark revision, scored population, and headline retrieval and answer-quality values.}

\textbf{Keywords}: Vaccine information, Vietnamese question answering, retrieval-augmented generation, hybrid retrieval, evidence attribution.

% ... sau \maketitle
\newpage
\begin{center}
    \bfseries\Large Acknowledgments
\end{center}

We sincerely thank all those who contributed to the completion of this capstone project.

First, our deepest appreciation goes to our supervisor, Mr. Tran Van Ha, whose expert guidance, thoughtful direction, and extensive domain knowledge were indispensable throughout every phase of this work. His constructive feedback, patience, and consistent encouragement enabled the team to refine both the conceptual foundations and technical implementation of this system to the highest standard we could achieve.

We are equally grateful to FPT University for providing a rigorous and stimulating academic environment. The knowledge imparted by faculty members across the department established the theoretical groundwork upon which this project was built. The university's infrastructure, research resources, and high standards of academic training created the conditions necessary for us to explore and realize our ideas with confidence.

We would also like to acknowledge our peers and colleagues for their candid discussions, willingness to share knowledge, and constructive critique. Their perspectives helped us examine the problem from multiple angles and materially improved the quality of this work.

Finally, and most profoundly, we thank our families for their unwavering support and understanding throughout this process. Their encouragement and the conditions they provided allowed us to devote ourselves fully to this research.
\newpage
\tableofcontents
\newpage
\listoffigures
\newpage
\section{Introduction}

\subsection{Motivation}
VaxiS addresses a practical information access problem. A user may know a vaccine by its product name and have a specific question, while the relevant source is a long instruction leaflet or regulatory document. Locating the applicable passage requires identifying the product, recognizing the relevant section and interpreting any qualifications recorded in nearby text or table notes. The project supports natural-language lookup of information in the collected vaccine usage instructions. The topic taxonomy in the older testdata module belongs to its dataset-generation workflow and does not define the current system's supported question categories.

Retrieval-augmented generation (RAG) brings evidence from an external collection into the generation process~\cite{lewis2020rag}. In this setting, the evidence should be accessible to the reader so the answer can be checked against the source. Producing a fluent response, citing a retrieved passage, and supporting each factual claim are distinct requirements. Research on generation with citations similarly evaluates attribution as a quality dimension alongside answer correctness~\cite{gao2023citations}.

The primary intended users are ordinary citizens. Healthcare workers and healthcare students are secondary users who may use the system for document lookup. The scope is general access to published information. Diagnosis, individual treatment decisions, and processing personal medical records fall outside the project scope.

\subsection{Objectives}
The main research question is how to retrieve relevant evidence from Vietnamese vaccine documents and use it to generate a comprehensible answer with traceable citations, including when the available evidence is incomplete.

\begin{enumerate}
\item Prepare a searchable corpus from available vaccine source documents, retaining the document identifiers and section metadata needed to trace evidence.
\item Compare lexical, dense, and hybrid retrieval, then assess the proposed Table-aware, Graph RAG, and Agent RAG extensions.
\item Generate Vietnamese responses that cite the supplied evidence and communicate when the evidence is insufficient.
\item Assess retrieval quality, answer quality, citation support, and abstention using documented datasets and reproducible run configurations.
\end{enumerate}

%
\section{Related Work}

\subsection{Automated Access to Vaccine Information}
Li et al.\ developed VaxBot-HPV around a curated knowledge base and a fine-tuned GPT model for HPV vaccine questions. Their evaluation included answer relevancy and faithfulness, two criteria that inform the assessment of VaxiS~\cite{li2025vaxbothpv}.

The Vaccine Information Resource Assistant (VIRA), developed by Johns Hopkins and IBM Research, uses natural language understanding to match questions to expert-vetted responses about COVID-19 vaccination. Its English and Spanish interfaces illustrate the value of organizing answers around user concerns and making reviewed content accessible in the user's language~\cite{jhuvira}.

K-Bot provides culturally tailored HPV vaccine information in Korean and English. Its authors designed dialogs around evidence-based content adapted to their intended audience. This work motivates attention to phrasing and accessibility in a Vietnamese-language interface~\cite{kim2025kbot}. VaxiS focuses on source-grounded question answering over the Vietnamese documents available to the project, with provenance carried through retrieval and generation.

\subsection{Evidence-Grounded Question Answering over Vaccine Documents}
BM25 uses corpus statistics and document length to score lexical matches~\cite{robertson2009bm25}. Dense retrieval instead ranks passages through similarity in an embedding space. BGE-M3 provides both dense representations and learned sparse lexical weights in a multilingual encoder~\cite{chen2024m3}. Reciprocal rank fusion combines the positions assigned by multiple retrieval streams on a common scale~\cite{cormack2009rrf}.

Evidence can also cross representational boundaries. HybridQA studies questions that require combining tables with associated passages~\cite{chen2020hybridqa}. Graph RAG research illustrates how entities and relationships extracted from documents can organize evidence for a query~\cite{edge2024graphrag}. ReAct illustrates how a language model can interleave decisions with actions that gather information~\cite{yao2023react}. These works motivate candidate extensions to the VaxiS pipeline. The specific designs and empirical effects of those extensions require project implementation records and completed runs.

\section{Methodology}

\subsection{Overview of the Vaccine Question-Answering System}
The documented pipeline has a document-preparation and indexing stage followed by question answering. Document preparation collects PDFs from the Drug Administration of Vietnam public service portal and converts them to Markdown using Google's Gemini. Indexing turns cleaned Markdown into identifiable chunks, adds document and section metadata, and persists the chunks in Parquet. The vector indexing path prepares dense and sparse representations in Qdrant. During question answering, the selected retriever ranks chunks. The generator receives the question and selected evidence, produces a Vietnamese response, and tracks the chunk identifiers cited by the response. The evaluation tooling also supports scoring ranked results stored from an earlier run.

\begin{figure}[htbp]
\centering
\fbox{\parbox{.9\linewidth}{\centering
Scraped PDFs $\rightarrow$ Gemini $\rightarrow$ Cleaned Markdown\\[.6em]
Section-aware chunks with metadata $\rightarrow$ Parquet and Qdrant\\[.6em]
Vietnamese question $\rightarrow$ Selected retriever $\rightarrow$ Evidence chunks\\[.6em]
Evidence and question $\rightarrow$ Generator $\rightarrow$ Response and source identifiers}}
\caption{Documented VaxiS indexing and question-answering flow. The evidence store used depends on the retrieval configuration.}
\label{fig:system}
\end{figure}

The repository includes a conversational entry point with query planning and memory, as well as single-query evaluation entry points. The benchmark protocol should record the entry point used. The presence of conversational planning in the application is insufficient to establish the Agent RAG experiment.

\subsection{Vaccine Metadata Module}
The corpus organization and vaccine name mapping associate cleaned source records with product identifiers and names. The implemented chunk schema supplies document identity, section headings, text, a topic category, a content type and a table flag. These fields support traceability and subsequent analysis, for example by identifying table-bearing chunks or grouping evidence by product. The document-discovery code concatenates the Markdown files in each registration folder into a single indexing document. It constructs a document ID from the mapped vaccine name and folder ID, and synthesizes a display title beginning with the Vietnamese phrase for vaccine usage instructions. This display title does not establish the title or type of each original PDF. The stored source-file metadata records Markdown filenames. The project plan also identifies document type, issuer, publication date and regulatory status as useful provenance fields.
\pending{Supply the frozen corpus metadata dictionary, field coverage, source URL and original-page mapping policy, and the rule for selecting among conflicting or superseded documents.}

\subsection{Vaccine Document Processing Module}
The team collected vaccine usage instructions from the Drug Administration of Vietnam public service portal at (\url{https://dichvucong.dav.gov.vn/congbothuoc/index}) on September 18, 2026. The listing was obtained with an HTTP GET request whose limit parameter was increased to retrieve all 123 vaccine listings in one request. The document extraction code supports several document fields, while the corpus selected for this project contains usage instructions. Some instructions also contain authorization-decision or regulatory-guidance information within the same PDF.

The preparation work was divided into three parts. PDFs were trimmed to retain the information needed for the project, removing label pages and other unnecessary material. Each file was uploaded individually to the Gemini website and converted to Markdown using the website's default model. Table images were also converted to HTML using Gemini, and the HTML was inserted into the Markdown files. The converted files were manually reviewed afterward. The exact Gemini model version was not recorded in the available materials.

The repository indexing code starts from the cleaned Markdown. Its parser splits the content into sections, and its chunker groups short sections and divides long sections under a length budget. The chunker defaults to a minimum of 80 and maximum of 600 estimated token units, with an overlap of 50. The estimate uses whitespace-separated words and an added allowance for HTML tags. The code accounts for Markdown and HTML table markup and handles oversized tables separately from prose. These are the current implementation defaults. \pending{Confirm the chunking settings used to build the corpus evaluated in the experiments.} The distinct Table-aware ablation still requires a versioned definition of its representation and retrieval behavior.

The design must preserve provenance when chunk boundaries are changed. A question can depend on a qualifier in an adjacent paragraph or a footnote that changes the interpretation of a row. A useful audit records the original source location and whether the selected evidence preserves those dependencies.
\pending{Provide the conversion prompt if retained, the conversion dates, and any recorded details of the manual review procedure. The exact Gemini model version remains unspecified unless a contemporaneous record is available.}

\subsection{Vaccine Knowledge Graph Module}
Graph RAG is the fifth configuration in the study. Its purpose is to test whether graph-structured evidence helps retrieve information distributed across passages. A reproducible specification requires the graph schema, the extraction and entity resolution rules, and the links from graph facts to source evidence.
\pending{Provide the Graph RAG implementation path and commit, node and edge schema, entity normalization and extraction procedure, storage backend, graph coverage and counts, and provenance strategy. Attach a verified graph example.}
\pending{Provide the graph-specific evaluation slice, comparison configuration, run artifacts, metrics and observed failure cases.}

\subsection{Retrieval-Augmented Generation Module}
Naive uses the repository's BM25 retrieval path over the chunk corpus. The evaluation CLI selects it with \artifact{--baseline}. Hybrid is the study baseline against which subsequent extensions are compared. Advanced embeds the original query with BGE-M3 and searches the dense Qdrant index before passing the top-ranked chunks to the shared generator.

Hybrid queries the BGE-M3 dense and learned sparse indexes. The implementation aggregates candidate ranks using weighted reciprocal rank fusion,
\[
 S(d)=\sum_{i:\,d\in L_i}\frac{w_i}{c+\mathrm{rank}_i(d)},
\]
where \(L_i\) is a candidate ranking, \(w_i\) its weight and \(c\) a smoothing constant. Hybrid uses the \artifact{BAAI/bge-m3} embedding checkpoint and the \artifact{BAAI/bge-reranker-v2-m3} cross-encoder reranker. The current configuration defaults to 50 candidates from each stream, a fusion constant of 60, 30 fused candidates, a final depth of five and enabled reranking. These are code defaults, and the experiment manifest must establish the settings actually used. The run record must identify the candidate depths, fusion settings, reranker checkpoint and final evidence depth.

Table-aware is the fourth configuration and is intended to measure retrieval and grounding on questions that depend on tabular structure. It keeps the Hybrid retrieval pipeline (BGE-M3 dense and sparse indexes, weighted reciprocal rank fusion and cross-encoder reranking) and changes how tabular evidence is indexed and surfaced. The shared chunker may split one table across several chunks, which separates rows from their column headers and footnotes. Table-aware therefore adds a preprocessing step before embedding and a merge step after fusion.

\paragraph{Index-time table preprocessing.}
Chunks flagged as tables are grouped when consecutive chunks of the same document share the same section header. Their HTML or Markdown fragments are merged into one table, and rows duplicated at chunk boundaries by the chunk overlap are removed. The identifiers of the merged chunks are stored in the chunk metadata so that evidence remains traceable to the original chunks.

Each merged table is then classified as \artifact{CLINICAL} or \artifact{STATISTICAL}. Clinical tables cover composition, dosage, schedule, adverse reactions, contraindications and storage. Statistical tables cover immunogenicity and efficacy results. The classifier checks keywords in the section header and table head, then the share of numeric cells (statistical at 0.45 or above), then the document category, and finally a 0.35 numeric-ratio threshold.

The language model \artifact{gpt-oss-120b} rewrites each table as Vietnamese prose, using a prompt chosen by table kind. Clinical tables are rewritten row by row, preserving every row, column, dose, frequency and footnote and naming the vaccine in each sentence. Statistical tables are summarized as one overview paragraph that names the study and control vaccines, age groups, pathogens or serotypes, and measures such as GMT, GMC and protective efficacy, without listing every number. Summaries are cached by a SHA-256 key over the document, header, table kind, model and merged table text. On a cache miss without LLM fallback, the merged table text is used alone.

Table-aware uses the \artifact{SUMMARIZE_DUAL} strategy by default. The prose summary and the original merged HTML table are stored together in a single chunk. The dense representation of this chunk benefits from the semantic summary, which improves matching with natural-language questions, while the sparse lexical representation can match exact terms and values from the original table. The generator receives the same chunk, so exact cell values remain available for answering.

\paragraph{Query-time table merging.}
The Table-aware configuration enables table chunk merging at query time. Retrieved table chunks that share a document and section header are merged after fusion and before reranking, and overlapping rows are removed. The merged chunk records its source chunk identifiers, so citations can still be mapped to the original chunks.

\pending{Provide the Table-aware implementation commit, the table-dependent query subset and its size, the run configuration, per-query outputs, retrieval and answer metrics, and an audited row or footnote error example.}

For Graph RAG, the retrieval component must specify how a question is associated with graph entities, how candidate paths are selected, and how their supporting passages enter the generator.
\pending{Provide the actual Graph RAG retrieval algorithm, graph search parameters, evidence ordering, and fallback behavior.}

Agent RAG is the sixth configuration and is intended to evaluate tool-mediated evidence gathering across retrieval steps.
\pending{Provide the Agent RAG implementation path and commit, model, available tools, routing and stopping logic, evidence aggregation and citation rules, call budget, fallback behavior, and an example trace.}
\pending{Provide the Agent RAG evaluation protocol, query coverage, comparison settings, metrics, tool-call and latency summaries, and failure analysis.}

\subsection{Answer Grounding and Verification}
The generator asks the model to answer from supplied evidence, cite chunk identifiers and communicate when the context is insufficient. The implementation extracts citation identifiers and checks their membership in the supplied chunks. It records valid and invalid identifiers and flags refusals by matching a fixed list of Vietnamese phrases. This heuristic requires validation against human judgments. The shared LLM response-validation function currently returns the provider response directly, so the code does not establish a general guarantee that every empty, refused or truncated completion is rejected. The planner applies additional output checks. Citation membership establishes traceability to context. Assessing whether a passage supports a particular claim requires semantic review.

The separate review application gives reviewers access to questions, reference answers, and source passages. Its workflow supports changes to the answer and evidence selection, with code paths for consensus and adjudication. The review application provides a way to audit the benchmark and its provenance.
\pending{Provide the final expert review protocol, reviewer qualifications and counts, adjudication outcomes, citation support rubric, and whether any separate automated claim-verification module contributed to the reported runs.}


\section{Dataset}
\label{sec:dataset}

\subsection{Source Corpus and Structural Distribution}
\label{subsec:corpus_distribution}
The documentary foundation of VaxiS is a curated corpus of 75 official Vietnamese vaccine usage-instruction leaflets (\emph{\vn{Tờ hướng dẫn sử dụng}}) collected from the Drug Administration of Vietnam (DAV) public service portal (\url{https://dichvucong.dav.gov.vn/congbothuoc/index}) on September 18, 2026. The initial registry query retrieved all 123 registered vaccine product listings. A systematic filtering protocol excluded 48 listings that lacked usable clinical text: 44 listings had no attached regulatory files, one failed due to a malformed download URL, two yielded empty directories, and one contained only packaging-label artwork without clinical usage instructions. The remaining 75 PDF documents correspond one-to-one with 75 canonical document identifiers (\artifact{doc_id}), each mapped to a unique vaccine product. In several instances, these official PDFs bundle marketing-authorization decisions or regulatory guidance alongside the clinical package insert.

Following PDF-to-Markdown conversion, HTML table reconstruction, and manual verification, the section parser and structure-aware chunker segment the 75 documents into 1,744 unique chunks. Each chunk is assigned a deterministic identifier formatted as \texttt{\{doc\_id\}\_c\_\{idx:04d\}}, alongside hierarchical section headers (\artifact{parent_header}, \artifact{header}), a content-type descriptor, a boolean table indicator (\artifact{metadata.has_table}), and a normalized medical category derived from 21 rule-based clinical heading patterns. By modality, 1,254 chunks (71.9\%) contain pure prose (\artifact{text}), 487 chunks (27.9\%) consist predominantly of tabular markup (\artifact{table}), and 3 chunks (0.2\%) combine prose with embedded tables (\artifact{mixed}), yielding 490 table-bearing chunks in total. Table~\ref{tab:corpus_categories} summarizes the corpus distribution across the standardized clinical categories.

\begin{table}[htbp]
\centering
\caption{Distribution of the 1,744 indexed corpus chunks across standardized medical categories.}
\label{tab:corpus_categories}
\begin{tabular}{llrr}
\toprule
Vietnamese Category Label & Clinical Domain (English) & Chunks & Share (\%)\\
\midrule
\vntt{dược lực học} & Pharmacodynamics and immunogenicity & 528 & 30.28\\
\vntt{tác dụng không mong muốn} & Adverse reactions and side effects & 224 & 12.84\\
\vntt{thông tin chung} & General product information & 141 & 8.08\\
\vntt{cảnh báo và thận trọng} & Warnings and clinical precautions & 140 & 8.03\\
\vntt{liều lượng và cách dùng} & Dosage, schedule, and administration & 102 & 5.85\\
\vntt{thành phần} & Qualitative and quantitative composition & 93 & 5.33\\
\vntt{bảng số liệu} & Standalone clinical and trial data tables & 74 & 4.24\\
\vntt{đặc tính dược lý} & Pharmacological properties & 55 & 3.15\\
\vntt{quá liều} & Overdose management & 55 & 3.15\\
\vntt{bảo quản} & Storage conditions and shelf life & 51 & 2.92\\
\vntt{chống chỉ định} & Contraindications & 51 & 2.92\\
\vntt{chỉ định} & Therapeutic and prophylactic indications & 45 & 2.58\\
\vntt{tương tác thuốc} & Drug and concomitant vaccine interactions & 44 & 2.52\\
\vntt{thai kỳ và cho con bú} & Pregnancy and lactation & 42 & 2.41\\
\multicolumn{2}{l}{\emph{Remaining 8 categories (packaging, pharmacokinetics, handling, etc.)}} & 99 & 5.68\\
\midrule
\textbf{Total} & & \textbf{1,744} & \textbf{100.00}\\
\bottomrule
\end{tabular}
\end{table}

As shown in Table~\ref{tab:corpus_categories}, pharmacodynamics and immunogenicity (\vntt{dược lực học}) and adverse reactions (\vntt{tác dụng không mong muốn}) account for over 43\% of all chunks. Because clinical trial outcomes and adverse-event frequencies are frequently formatted as multi-column tables that exceed the 600-token chunking budget, the baseline chunker splits large tables across consecutive sub-chunks while repeating column headers. In the raw indexing configuration, all 1,744 chunks are indexed directly into the baseline vector collection. In contrast, the table-aware indexing pipeline merges consecutive split table chunks sharing the same \texttt{(doc\_id, header)} pair and deduplicates overlapping boundary rows. This consolidation collapses the 490 table-bearing sub-chunks into 212 complete tables—classified automatically into 143 clinical guidance tables and 69 statistical immunogenicity tables—reducing the indexed vector collection from 1,744 to 1,466 points while retaining all constituent sub-chunk identifiers in chunk metadata for provenance alignment.

\subsection{Stratified 300-Query Benchmark Construction}
\label{subsec:benchmark_construction}
To evaluate retrieval and grounded generation across diverse information-seeking scenarios, we constructed a 300-question Vietnamese benchmark. The benchmark was synthesized from the chunked corpus using DeepEval's context-grounded \artifact{Synthesizer} (one golden per context group, with expected outputs)~\cite{deepevalcontexts}, driven by \artifact{gpt-oss-120b} (temperature $0.2$, deterministic random seed $42$, generated on November 16, 2026).

Rather than sampling chunks uniformly—which would over-represent dense trial tables—the generator enforces a three-dimensional stratified design across three orthogonal axes:
\begin{enumerate}
\item \textbf{Answerability (70\% Answerable, 30\% Partially Answerable):} \emph{Answerable} queries ($n=210$) ask for clinical facts fully supported by the supplied context group. \emph{Partially Answerable} queries ($n=90$) are deliberately formulated as two-part inquiries: one clause is directly answerable from the reference context, whereas the second clause introduces an out-of-context clinical constraint or secondary question, testing whether the system can answer the supported component while explicitly abstaining from the unsupported portion.
\item \textbf{Reasoning Complexity (50\% Single-Hop, 30\% Multi-Hop, 20\% Table Reasoning):} \emph{Single-Hop} items ($n=150$) are synthesized from a single non-table prose chunk using \emph{Concretizing} ($0.5$) and \emph{Constrained} ($0.5$) query evolutions. \emph{Table Reasoning} items ($n=60$) target a single table-bearing chunk (\artifact{has_table=True}) under \emph{Reasoning} ($0.5$) and \emph{Comparative} ($0.5$) evolutions, requiring precise lookup or comparison of dosage intervals, age thresholds, seroprotection rates, or geometric mean titers (GMT). \emph{Multi-Hop} items ($n=90$) require synthesizing evidence across two distinct chunks and are balanced into an exact 50/50 split between two sub-types:
\begin{itemize}
\item \emph{Intra-document multi-hop (\artifact{same_doc}, $n=45$):} Pairs two clinically complementary sections from the same vaccine leaflet across 11 predefined category pairs (e.g., \emph{Indications $\times$ Contraindications}, \emph{Dosage $\times$ Warnings}, \emph{Composition $\times$ Contraindications}), excluding uninformative overdose stubs under 200 characters, and evolves queries via \emph{Multicontext} ($0.7$) and \emph{Reasoning} ($0.3$) transformations.
\item \emph{Inter-vaccine comparative multi-hop (\artifact{cross_vaccine}, $n=45$):} Pairs chunks belonging to the same clinical category across two different vaccine products (spanning eight comparable domains such as storage, pregnancy, adverse reactions, or schedules) under \emph{Comparative} ($0.6$) and \emph{Multicontext} ($0.4$) evolutions, requiring explicit naming of both vaccines.
\end{itemize}
\item \textbf{Query Style (60\% Patient Style, 40\% Clinical Style):} \emph{Patient Style} queries ($n=180$) emulate layperson inquiries on health forums or messaging platforms, featuring everyday Vietnamese phrasing, conversational particles, and brief personal context. \emph{Clinical Style} queries ($n=120$) emulate healthcare professionals verifying immunization protocols for complex patient cases with multiple clinical constraints (e.g., comorbidities, delayed booster intervals, concomitant medications).
\end{enumerate}

The joint allocation across the $2 \times 3 \times 2 = 12$ cells is computed via the Hamilton largest-remainder method followed by iterative marginal reconciliation, ensuring that the cell counts sum to 300 while preserving the exact target marginals along every axis (Table~\ref{tab:dataset}). Table~\ref{tab:benchmark_examples} illustrates representative queries across these strata.

\begin{table}[htbp]
\centering
\caption{Joint three-axis stratified distribution and multi-rater expert review allocation across the 12 benchmark strata ($N=300$ queries, 435 total review turns across 3 clinical raters).}
\label{tab:dataset}
\begin{tabular}{lllrrrrr}
\toprule
Answerability & Complexity & Query Style & Total & Calib.\ & Overlap & Indep.\ & Audit\\
 & & & ($N$) & (3 raters) & (3 raters) & (per rater) & (per rater)\\
\midrule
Answerable & Single-hop & Patient & 63 & 3 & 6 & $18 \times 3$ & $0 \times 3$\\
(210; 70\%) & Single-hop & Clinical & 42 & 2 & 4 & $12 \times 3$ & $0 \times 3$\\
 & Multi-hop$^\dagger$ & Patient & 38 & 2 & 3 & $11 \times 3$ & $0 \times 3$\\
 & Multi-hop$^\dagger$ & Clinical & 25 & 1 & 3 & $7 \times 3$ & $0 \times 3$\\
 & Table reasoning & Patient & 25 & 1 & 3 & $7 \times 3$ & $3 \times 3$\\
 & Table reasoning & Clinical & 17 & 1 & 1 & $5 \times 3$ & $2 \times 3$\\
\midrule
Partially & Single-hop & Patient & 27 & 1 & 2 & $8 \times 3$ & $2 \times 3$\\
Answerable & Single-hop & Clinical & 18 & 1 & 2 & $5 \times 3$ & $1 \times 3$\\
(90; 30\%) & Multi-hop$^\dagger$ & Patient & 16 & 1 & 3 & $4 \times 3$ & $1 \times 3$\\
 & Multi-hop$^\dagger$ & Clinical & 11 & 1 & 1 & $3 \times 3$ & $1 \times 3$\\
 & Table reasoning & Patient & 11 & 1 & 1 & $3 \times 3$ & $3 \times 3$\\
 & Table reasoning & Clinical & 7 & 0 & 1 & $2 \times 3$ & $2 \times 3$\\
\midrule
\textbf{Total Questions} & & & \textbf{300} & \textbf{15} & \textbf{30} & \textbf{255} & \textbf{45}\\
\textbf{Review Turns} & & & \textbf{435} & \textbf{45} & \textbf{90} & \textbf{255} & \textbf{45}\\
\bottomrule
\multicolumn{8}{l}{\footnotesize $^\dagger$Across the 90 Multi-hop queries, 45 are intra-document (\artifact{same_doc}) and 45 are inter-vaccine (\artifact{cross_vaccine}).}
\end{tabular}
\end{table}

\begin{table}[htbp]
\centering
\caption{Representative queries from the benchmark illustrating the stratification dimensions.}
\label{tab:benchmark_examples}
\begin{tabularx}{\linewidth}{llX}
\toprule
Query ID & Stratum & Vietnamese Query (\artifact{input}) and Ground-Truth Provenance\\
\midrule
\artifact{eval\_0001} & Answerable / Single-hop / Patient & \emph{\vn{``Phụ nữ 22 tuần mang thai có nên tiêm vaccine Abrysvo không ạ?''}} (Source: \artifact{abrysvo\_540310322125\_c\_0006}; gestational window 24--36 weeks.)\\
\artifact{eval\_0125} & Answerable / Multi-hop (\artifact{cross\_vaccine}) / Patient & \emph{\vn{``Vắc-xin Peg-Grafeel và Boostrix khác nhau như thế nào về hạn sử dụng và cách bảo quản?''}} (Sources: \artifact{peg\_grafeel\_qlsp\_0636\_13\_c\_0006}, \artifact{boostrix\_300310036923\_c\_0014}.)\\
\artifact{eval\_0194} & Answerable / Table reasoning / Clinical & \emph{\vn{``Trong nhóm tuổi 3--17 năm, so sánh mức trung bình của tiêu đề kháng thể (average titers) giữa vaccine Influvac Tetra và vaccine Influvac 1/2 đối với các chủng virus A/H1N1, B(Yam) và B(Vic) là như thế nào?''}} (Source: \artifact{influvac\_tetra\_870310304024\_c\_0005}.)\\
\artifact{eval\_0238} & Partially Answerable / Single-hop / Clinical & \emph{\vn{``Trong bệnh nhân 68 tuổi, có tiền sử bệnh tim mạch và đã tiêm liều đầu Shingrix, liệu có thể đồng thời dùng aspirin 81 mg hàng ngày và tiêm liều thứ hai của Shingrix theo lịch trễ mũi 2 tháng không? Ngoài ra, liệu có thể tiêm đồng thời vaccine COVID-19 mRNA cùng lúc với Shingrix không?''}} (Source: \artifact{shingrix\_540310303224\_c\_0007}.)\\
\artifact{VAX\_Q\_0001} & Unanswerable / Auxiliary refusal set & \emph{\vn{``Vaxneuvance có gây sưng đỏ đuôi mắt không?''}} (Source: none; the leaflet does not document this reaction, so the expected behavior is abstention.)\\
\bottomrule
\end{tabularx}
\end{table}

Because DeepEval's asynchronous synthesis returns generated items out of input batch order, naive positional indexing in the initial export caused misalignment between each generated context and its recorded source chunk identifiers in 279 of the 300 records. To guarantee provenance integrity, we implemented an offline deterministic context-to-chunk reconciliation pass. This procedure matches each embedded context string against the canonical \texttt{[doc\_title | category | header]} formatting of the chunk store and enforces structural invariants (verifying that \artifact{same_doc} items reference exactly one \artifact{doc_id}, \artifact{cross_vaccine} items reference two distinct \artifact{doc_ids}, and \artifact{TABLE_REASONING} items include at least one chunk with \artifact{has_table=True}). The resulting benchmark achieves 100\% provenance alignment with the corpus and serves as the ground-truth benchmark for all retrieval and generation experiments.

For stress-testing complete abstention on wholly unanswerable queries, we construct a dedicated auxiliary set of unanswerable questions. The seed pool consists of 1,200 Vietnamese consumer health questions released in Stage~3 of the R2AI 2026 competition, most of which are not about vaccines. Each seed question is rewritten by Qwen2.5-7B-Instruct. The model receives the list of vaccines available in our corpus and must produce a natural Vietnamese question that names exactly one of these vaccines while preserving the clinical intent and any symptoms or conditions of the original query. Because the original concerns are largely unrelated to vaccination, most rewritten questions cannot be answered from the corresponding leaflet. From the resulting pool we retain 100 questions labeled \emph{Unanswerable}. These carry no gold source chunks, since the expected behavior is abstention rather than retrieval.

\subsection{Expert Review and Quality Assurance Protocol}
\label{subsec:expert_review_protocol}
To audit label validity, reference-answer accuracy, and evidence completeness, we developed a dedicated web-based clinical review platform backed by a stratified three-expert allocation protocol. As detailed in the right-hand columns of Table~\ref{tab:dataset}, the 300 benchmark questions are partitioned deterministically (seed $42$) into four functional batches across three medical reviewers:
\begin{enumerate}
\item \textbf{Calibration Batch (15 questions, 45 review turns):} A stratified 5\% sample evaluated independently by all three experts prior to consensus discussion, aligning interpretation of the annotation rubric and partial-answerability boundaries.
\item \textbf{Inter-Annotator Overlap Batch (30 questions, 90 review turns):} A stratified 10\% sample scored independently and blindly by all three reviewers to quantify inter-annotator agreement (IAA). The platform computes Fleiss' Kappa ($\kappa$)~\cite{fleiss1971measuring} with Landis and Koch agreement tiers~\cite{landis1977measurement} across two dimensions—three-class answerability classification and binary chunk-level evidence selection—supplemented by mean pairwise Jaccard similarity across the selected evidence sets.
\item \textbf{Independent Primary Batch (255 questions, 85 per reviewer):} The remaining 85\% of the benchmark is divided equally into three balanced, non-overlapping packages of 85 questions per expert, preserving identical marginal proportions across all 12 strata.
\item \textbf{Risk-Prioritized Cross-Audit Batch (45 questions, 15 per reviewer):} To safeguard high-complexity items within the independent packages, 15 questions from each expert's primary allocation (45 questions in total) are routed in a round-robin schedule (each expert's items are audited by the next expert in cyclic order) to a second expert for cross-auditing. The risk-based selection rule concentrates 100\% of these cross-audit slots on the most error-prone strata—specifically \emph{Table Reasoning} (30 audited items) and \emph{Partially Answerable} non-table items (15 audited items)—yielding exactly 145 total evaluation turns per reviewer (435 turns overall).
\end{enumerate}

Within the review interface, experts inspect the generated question, expected answer, and highlighted source chunks alongside a searchable full-document view. Reviewers can approve an item as-is (\artifact{APPROVED}), modify the question, reference answer, answerability label, or supporting chunk list with structured rationales (\artifact{EDITED}), or flag malformed synthetic items (\artifact{REJECTED}), thereby establishing an auditable provenance chain from raw regulatory PDFs to benchmark evaluation.


\section{Experiments and Results}
\label{sec:experiments}

\subsection{Experimental Design}
The study compares the six configurations in Table~\ref{tab:sequence}. Each comparison should use a frozen corpus and judgment set, the same scored questions and matched evidence budgets. Record differences intrinsic to the methods, such as using a graph traversal or multiple retrieval calls. Hold the generator, prompts and judge settings fixed for comparisons intended to isolate the effect of retrieval. Hybrid serves as the reference for extensions. Because Hybrid can change both retrieval streams and reranking, the sequence is a system-level comparison. Component effects require matched controls, such as Hybrid with reranking disabled.

\begin{table}[htbp]
\centering
\caption{Ablation-study sequence and documented evidence paths.}
\label{tab:sequence}
\begin{tabularx}{\linewidth}{rlX}
\toprule
Order & Configuration & Evidence path\\
\midrule
1 & Naive & BM25 over corpus chunks.\\
2 & Advanced & One BGE-M3 dense query against Qdrant.\\
3 & Hybrid (baseline) & Dense and learned sparse retrieval, rank fusion and optional reranking.\\
4 & Table-aware & \pending{Record final table representation and evidence path.}\\
5 & Graph RAG & \pending{Record final graph construction and retrieval path.}\\
6 & Agent RAG & \pending{Record final orchestration and tool-use path.}\\
\bottomrule
\end{tabularx}
\end{table}

For retrieval, Hit@$k$ indicates whether at least one labeled relevant chunk was retrieved by rank \(k\). Recall@$k$ measures the proportion of labeled evidence recovered and Precision@$k$ measures the proportion of the first \(k\) positions occupied by labeled evidence. Mean reciprocal rank rewards early retrieval of the first labeled source. Normalized discounted cumulative gain rewards placing labeled sources near the top. Alternative supporting chunks can receive no credit when the gold evidence is incomplete. Report scoring coverage and the number of rankings shorter than each cutoff.

The evaluation code also supports mean average precision at a cutoff and LLM-judged answer relevancy and faithfulness, contextual relevancy, precision and recall, and a default-enabled GEval Correctness metric against the reference answer. Citation analysis should distinguish citation validity from claim support, and failure analysis should assess whether refusals and partial answers match the evidence available. The source ODT identifies user testing as an evaluation objective.
\pending{Provide the final run manifest with code revision, index revision, checkpoint identifiers, prompts, judge settings, hardware, retrieval and context budgets, inclusion and exclusion rules, and significance or uncertainty procedure if used. Supply the human review and user-testing protocols and completed results if available.}

\subsection{Results}
The saved generation artifact holds 300 records, each with four retrieved identifiers. The recorded generation model is \artifact{gpt-oss-120b}. All labeled source IDs in the corrected benchmark and all retrieved IDs in these saved outputs occur in the current Parquet corpus. Identifier coverage alone leaves the corpus and index revision used for the saved run unresolved. The saved latency field is measured within answer generation and excludes retrieval time. This artifact provides evidence for scoring those saved rankings at cutoffs up to four. A score calculated above that depth would have missing positions and would be unsuitable as evidence for a deeper live retrieval experiment. Run metadata and comparable per-query outputs are required to attribute a result to an ablation configuration.

\begin{table}[htbp]
\centering
\caption{Evidence required to complete the ablation results.}
\label{tab:results}
\begin{tabularx}{\linewidth}{lX}
\toprule
Configuration & Pending results\\
\midrule
Naive & \pending{Insert scored query count, retrieval metrics, answer and citation scores, abstention and latency from the BM25 run.}\\
Advanced & \pending{Insert matched dense retrieval, answer and citation scores, abstention and latency.}\\
Hybrid (baseline) & \pending{Identify the baseline run and insert matched retrieval, answer and citation scores, abstention and latency.}\\
Table-aware & \pending{Insert verified full-set and table-slice metrics and deltas from Hybrid, including row or cell attribution and latency.}\\
Graph RAG & \pending{Insert verified full-set and multi-hop metrics and deltas from Hybrid, including graph provenance and latency.}\\
Agent RAG & \pending{Insert verified matched metrics and deltas from Hybrid, with call counts, failures and latency.}\\
\bottomrule
\end{tabularx}
\end{table}

\pending{Add the final numerical results table with denominators and cutoffs and interpret only observed differences. Include per-slice results for single-hop, multi-hop, table reasoning and partial answerability, supported by saved run artifacts.}

\subsection{Qualitative Assessment of Vietnamese Vaccine Answers}
A qualitative assessment should trace the question to the ranked evidence and examine the generated response claim by claim. For product-specific questions, check whether the passage refers to the right product and formulation. For tables, check whether the evidence includes the relevant headers and any qualifying notes. For evidence distributed across sources, check whether all required passages appear and whether the answer combines them consistently. A partially answerable case should show which parts the source supports and how the response handles the remaining uncertainty.

\pending{Insert reviewed examples with query IDs, variant, source titles and chunk IDs, source excerpts, actual output and assessor judgments. Include a supported answer, an evidence miss or unsupported claim, a table-dependent example and a partially answerable case. Supply reviewed Graph RAG and Agent RAG examples once the corresponding runs are available.}

\section{Conclusion}
\label{sec:conclusion}
VaxiS applies evidence-grounded retrieval and generation to Vietnamese vaccine information. The available project materials establish an indexing path from cleaned Markdown, BM25, dense and hybrid retrieval implementations, chunk-level citation tracking, a working benchmark, and tooling for reviewing reference answers and evidence. These artifacts support evaluation of whether relevant passages are retrieved and whether the resulting answers can be checked against those passages.

The central empirical question is how the six configurations compare under consistent evaluation conditions. The documented corpus and benchmark audit, a reproducible comparative run, and verification of the final three configurations are prerequisites for a comparative conclusion.
\pending{Summarize the verified findings, cite the completed results tables, state observed failure modes and delimit the conclusions to the evaluated corpus and questions.}

\bibliographystyle{unsrtnat}
\bibliography{references}

% --- Bibliography ---
\bibliographystyle{unsrtnat}
\bibliography{references}

% --- Appendix ---
\appendix
\section*{Appendix}
\addcontentsline{toc}{section}{Appendix}
\label{sec:appendix}
\setcounter{subsection}{0}
\renewcommand{\thesubsection}{\Alph{subsection}}
\renewcommand{\theHsubsection}{appendix.\Alph{subsection}}

% ------------------------------------------------------------------
% Appendix material: prompts, schemas, and reproducibility configs
% (Pseudocode stages intentionally omitted.)
% ------------------------------------------------------------------
%\lstset{
%    basicstyle=\ttfamily\footnotesize,
%    breaklines=true,
%    columns=fullflexible,
%    frame=single,
%    rulecolor=\color{black!20},
%    showstringspaces=falsface
\subsection{Source-to-Markdown Processing, Metadata Extraction, and Chunking}
Collection is documented in \artifact{scrape_dataset/scrape_data.py}, \artifact{scrape_dataset/extract_documents.py}, and \artifact{scrape_dataset/download_documents.py}. The team used the Gemini website with its default model, uploading one trimmed PDF at a time. Table images were converted to HTML and inserted into the Markdown, followed by manual review. \pending{Provide the conversion prompt or retained session record if available.}
The principal indexing files are \artifact{src/vacxis_capstone_fa26/indexing/parser.py}, \artifact{src/vacxis_capstone_fa26/indexing/chunker.py}, and \artifact{src/vacxis_capstone_fa26/indexing/indexer.py}. Review the original source alongside its cleaned Markdown and indexed output when checking the preservation of a table or footnote.
\pending{Provide a representative source-to-Markdown example, Gemini conversion and cleaning settings, and the indexing command and resulting chunk metadata for the final corpus.}

\subsection{Answer Generation Prompts and Output Schema}
The generation files \artifact{src/vacxis_capstone_fa26/generation/prompts.py} and \artifact{src/vacxis_capstone_fa26/generation/generator.py} define the prompt construction and citation handling. Saved result records include the answer, valid and invalid citations, cited and retrieved IDs, a refusal flag, model, latency and usage.
\pending{Insert the exact versioned generator prompt, context formatting, model settings and a sanitized saved response record from the final comparison.}

\subsection{Graph Schema and Constraints}
\pending{Insert the versioned Graph RAG node and edge specification, graph serialization or storage constraints, provenance fields, and a worked retrieval trace with source identifiers.}

\subsection{Example Questions and Reviewed Outputs}
\pending{Insert approved benchmark examples with expected and actual outputs, evidence IDs, and reviewer notes. Include each evaluated configuration where comparable examples are available.}
\subsection{PDF-to-Markdown Processing, GRI-Aware Extraction, and Chunking}

\paragraph{Pipeline configuration (paths and model handle):}
\begin{lstlisting}[language=Python,caption={Module 1 pipeline configuration fields (paths and model handle).},label={lst:m1-pipeline-config}]
PipelineConfig:
    input_pdf: Optional[Path]
    output_dir: Path = "./data"
    llm_model: str = ".../Qwen3-4B-Instruct-2507"
    gri_taxonomy_path: Path = "data/taxonomies/gri_taxonomies_by_code.json"

Derived outputs:
    split_dir:   output_dir/"1_split"/<report_name>/
    markdown_dir:output_dir/"2_markdown"/<report_name>/
    merged_file: output_dir/"3_merged"/<report_name>.md
    output_json: output_dir/"4_output"/<report_name>.json
\end{lstlisting}

\paragraph{TOC extraction prompt (JSON array output):}
\begin{lstlisting}[language=json,caption={Module 1 TOC extraction prompt (JSON array).},label={lst:m1-toc-prompt}]
Extract the Table of Contents from the document content.

Return ONLY a JSON array with the following format:
[
    {"title": "Section title", "page": "page number or null", "level": 1},
    ...
]

Rules:
- level 1 for main sections, level 2 for subsections
- If no page number, use null
- Extract ALL section titles from the table of contents

Document content:
<pages_content>
\end{lstlisting}

\paragraph{GRI reference table extraction prompt (JSON array output):}
\begin{lstlisting}[language=json,caption={Module 1 GRI reference extraction prompt (JSON array).},label={lst:m1-gri-ref-prompt}]
Extract ALL GRI references from the document content.

Return ONLY a JSON array with the following format:
[
    {"title": "Disclosure title", "indicator": "GRI X-X", "page": "page number or null"},
    ...
]

Rules:
- indicator format: "GRI X-X" (e.g., "GRI 2-2", "GRI 302-1")
- Extract ALL GRI indicators found in the document

Document content:
<pages_content>
\end{lstlisting}

\paragraph{Page-level extraction output schema (JSON):}
\begin{lstlisting}[language=json,caption={Module 1 page-level extraction schema (sections + provenance).},label={lst:m1-page-extraction-schema}]
{
    "has_gri_content": true,
    "sections": [
        {
            "title": "Section title",
            "content": "Content",
            "gri_indicator": "GRI X-X",
            "gri_source": "content|context|reference",
            "is_continuation": false
        }
    ],
    "summary": "Brief summary"
}
\end{lstlisting}

\paragraph{End-to-end JSON output structure and chunking metadata:}
\begin{lstlisting}[language=json,caption={Module 1 output JSON top-level keys and chunking metadata.},label={lst:m1-output-json-schema}]
{
    "table_of_contents": [
        {"title": "...", "page": "...", "level": 1}
    ],
    "gri_references": [
        {"title": "...", "indicator": "GRI X-X", "page": "..."}
    ],
    "page_extractions": [
        {"page": 1, "sections": [ ... ], "summary": "..."}
    ],
    "chunking_metadata": {
        "max_tokens": 512,
        "chunk_size": 200,
        "chunk_overlap": 50,
        "total_original_sections": 0,
        "total_chunked_sections": 0
    }
}
\end{lstlisting}

\subsection{Entity Extraction Prompts and Output Schema}

\paragraph{System prompt:}
\begin{lstlisting}[language=json,caption={Module 2 system prompt (JSON-only output constraint).},label={lst:m2-system-prompt}]
You are an expert in ESG (Environmental, Social, Governance) reporting and GRI (Global Reporting Initiative) standards.
Your task is to extract structured information from ESG report text and align it with relevant GRI disclosure indicators.

You must return valid JSON only, without any markdown formatting or explanations.
\end{lstlisting}

\paragraph{Single-chunk extraction prompt (template):}
\begin{lstlisting}[caption={Module 2 single-chunk extraction prompt (template; placeholders shown).},label={lst:m2-single-prompt}]
Given the following text chunk from an ESG report, extract structured information.

Report Context:
- Company: {company_name}
- Report Type: {report_type}
- Report Year: {report_year}
- ESG Category: {esg_category} ({esg_category_full})

Text Chunk:
{chunk_text}

Relevant GRI Standards for {esg_category} category:
{gri_context}

Task:
1) ESG Actions: name, description
2) Key Statements: text, importance (high/medium/low)
3) GRI Alignments: list GRI codes from the provided list above

Output Format (JSON only):
{
    "actions": [{"name": "...", "description": "..."}],
    "statements": [{"text": "...", "importance": "high"}],
    "gri_alignments": ["GRI XXX-Y"]
}

Return only valid JSON without any markdown formatting.
\end{lstlisting}

\paragraph{Batch extraction prompt (schema constraints):}
\begin{lstlisting}[caption={Module 2 batch extraction prompt (schema constraints; shortened).},label={lst:m2-batch-prompt}]
CRITICAL INSTRUCTIONS:
1. Return ONLY a valid JSON array. No introduction, no markdown.
2. The array MUST have exactly {chunk_count} objects, one per input chunk in the same order.
3. Each object MUST have exactly these keys: chunk_id, actions, statements, gri_alignments
4. If a chunk has no ESG content, return empty arrays.

OUTPUT SCHEMA:
[
    {
        "chunk_id": <number>,
        "actions": [{"name": "<short>", "description": "<1-2 sentences>"}],
        "statements": [{"text": "<verbatim>", "importance": "high|medium|low"}],
        "gri_alignments": ["GRI XXX-Y"]
    }
]
\end{lstlisting}

\subsection{Graph JSON Schema and Neo4j Constraints}

\paragraph{Graph JSON structure:}
\begin{lstlisting}[language=json,caption={Module 2 graph JSON top-level structure (schema).},label={lst:m2-graph-json-schema}]
{
    "metadata": {
        "company": "<string>",
        "report_type": "<string>",
        "report_year": <int>,
        "node_count": <int>,
        "relationship_count": <int>
    },
    "nodes": [
        {"label": "Company|Report|ESG_Category|ESG_Action|ESG_Statement|GRI_Disclosure",
         "properties": {"...": "..."}}
    ],
    "relationships": [
        {"type": "HAS_REPORT|HAS_STATEMENT|BELONGS_TO_CATEGORY|CATEGORIZED_AS|BELONGS_TO_ACTION|ALIGNS_TO",
         "from_label": "<label>", "from_id": "<id>",
         "to_label": "<label>", "to_id": "<id>",
         "properties": {"...": "..."}}
    ]
}
\end{lstlisting}

\paragraph{Uniqueness constraints (Neo4j):}
Table~\ref{tab:edges} summarizes the relationship types used in the knowledge graph. The constraints below focus only on node-level uniqueness keys used to prevent duplicate entity creation.
\begin{lstlisting}[caption={Module 2 Neo4j uniqueness constraints (conceptual list).},label={lst:m2-neo4j-constraints}]
Company(company_id)
Report(report_id)
ESG_Statement(statement_id)
ESG_Action(action_id)
ESG_Category(category_id)
GRI_Disclosure(disclosure_id)
\end{lstlisting}

To complement the prompt and schema definitions above, the following examples illustrate how the system handles real benchmark questions from retrieval to answer verification.

\subsection{Example Questions and Fact-Checked Outputs}
\label{sec:appendix_examples}

This section provides three representative examples of the question-answering and fact-checking workflow across Environmental (E), Social (S), and Governance (G) topics. Each example includes the input question, generated answer, retrieved report evidence, and claim-level fact-checking results. The fact-checking examples include both supported and non-supported claims to illustrate how the verifier detects partially unsupported or incorrect statements.

%% =================== Box 1 — E (ACB GHG emissions) ===================
\begin{tcolorbox}[colback=white,colframe=black!65,boxrule=0.6pt,arc=0pt,left=8pt,right=8pt,top=8pt,bottom=8pt]

\textbf{Question:} Did ACB reduce its total GHG emissions across all three scopes in 2023 compared to 2022?

\textbf{Generated answer:} Yes, ACB reduced its total GHG emissions across all three scopes in 2023 compared to 2022. The total GHG emissions decreased by 9\%, from 34,355 tonnes of CO$_2$ in 2022 to 31,238 tonnes of CO$_2$ in 2023 (ACB, 2023, Page 34). This reduction is also supported by GRI Standard GRI 305-1, which pertains to direct GHG emissions.

\begin{itemize}
    \item \textbf{Claim:} ACB's total GHG emissions were 31,238 tonnes of CO$_2$ in 2023.
    \begin{itemize}
        \item \textbf{Label:} SUPPORTS
        \item \textbf{Confidence:} 0.90
        \item \textbf{Source:} \href{https://static2.vietstock.vn/vietstock/2025/3/26/20250325_acb_250325_annual_report_2025.pdf}{vietstock.vn (ACB Annual Report)}
    \end{itemize}

    \item \textbf{Claim:} ACB's total GHG emissions decreased by 9\% in 2023 compared to 2022.
    \begin{itemize}
        \item \textbf{Label:} NOT\_ENOUGH\_INFO
        \item \textbf{Confidence:} 0.30
        \item \textbf{Source:} \href{https://climate.ec.europa.eu/news-other-reads/news/climate-action-progress-report-2023-shows-largest-annual-drop-emissions-decades-2024-10-31_en}{climate.ec.europa.eu (EU Climate Action Progress Report 2023)}
    \end{itemize}
\end{itemize}

\end{tcolorbox}


% =================== Box 2 — S (MSB E&S risk loan assessment) ===================
\begin{tcolorbox}[colback=white,colframe=black!65,boxrule=0.6pt,arc=0pt,left=8pt,right=8pt,top=8pt,bottom=8pt]

\textbf{Question:} How many loans were assessed for Environmental and Social risks by MSB in 2024?

\textbf{Generated answer:} 84,619. In 2024, MSB conducted Environmental and Social risk assessments for 84,619 loans out of a total of 134,487 loan proposal documents (MSB, 2024, Page 20). This assessment aligns with GRI Standard GRI 2-21, which pertains to the management of environmental and social risks.

\begin{itemize}
    \item \textbf{Claim:} In 2024, MSB conducted Environmental and Social risk assessments for 84,619 loans out of a total of 134,487 loan proposal documents.
    \begin{itemize}
        \item \textbf{Label:} NOT\_ENOUGH\_INFO
        \item \textbf{Confidence:} 0.95
        \item \textbf{Source:} \href{https://cmsv5.fiingroup.vn/medialib/FG/2025/2025-04/2025-04-22/20250422_MSB-250422-Annual-Report-2024.pdf}{fiingroup.vn (MSB Annual Report 2024)}
    \end{itemize}
\end{itemize}

\end{tcolorbox}


% =================== Box 3 — G (HDBank TCFD participation) ===================
\begin{tcolorbox}[colback=white,colframe=black!65,boxrule=0.6pt,arc=0pt,left=8pt,right=8pt,top=8pt,bottom=8pt]

\textbf{Question:} Did HDBank participate in LeapFrog's first TCFD report in 2023?

\textbf{Generated answer:} Yes, HDBank participated in LeapFrog's first Task Force on Climate-related Financial Disclosures (TCFD) report in 2023. This involvement is significant for assessing climate change risks and opportunities within LeapFrog's investment portfolio (HDBank, 2023, Page 12; GRI: GRI 405-1, GRI 202-1, GRI 201-2).

\begin{itemize}
    \item \textbf{Claim:} HDBank participated in LeapFrog's first TCFD report in 2023.
    \begin{itemize}
        \item \textbf{Label:} SUPPORTS
        \item \textbf{Confidence:} 0.80
        \item \textbf{Source:} \href{https://cdn.hdbank.com.vn/hdbank-file/investor/src/bctn-2023/phat-trien-ben-vung-EN.html}{hdbank.com.vn (Sustainability Report 2023)}
    \end{itemize}
\end{itemize}

\end{tcolorbox}


% =================== Box 4 — E (BIDV green credit 2023) ===================
\begin{tcolorbox}[colback=white,colframe=black!65,boxrule=0.6pt,arc=0pt,left=8pt,right=8pt,top=8pt,bottom=8pt]

\textbf{Question:} What was the total outstanding balance of green credit at BIDV as of December 31, 2023 (in VND billion)?

\textbf{Generated answer:} 74,177 VND billion. As of December 31, 2023, BIDV reported a total outstanding balance of green credit amounting to 74,177 VND billion, which accounted for 4.24\% of its total outstanding loans (BIDV, 2023, Page 37). This information is also supported by GRI Standard GRI 308-1, which pertains to the organization's green credit performance.

\begin{itemize}
    \item \textbf{Claim:} As of December 31, 2023, BIDV reported a total outstanding balance of green credit amounting to 74,177 VND billion.
    \begin{itemize}
        \item \textbf{Label:} SUPPORTS
        \item \textbf{Confidence:} 0.95
        \item \textbf{Source:} \href{https://www.studocu.vn/vn/document/hoc-vien-tai-chinh/tai-chinh-tien-te/bidv-annual-report-2023-fn-da-nen-da-nen-1/108663677}{studocu.vn (BIDV Annual Report 2023)}
    \end{itemize}

    \item \textbf{Claim:} As of December 31, 2023, BIDV's green credit accounted for 4.24\% of its total outstanding loans.
    \begin{itemize}
        \item \textbf{Label:} NOT\_ENOUGH\_INFO
        \item \textbf{Confidence:} 0.70
        \item \textbf{Source:} \href{https://bidv.com.vn/bidv_en/tin-tuc/tin-ve-bidv/bidv-danh-10000-ty-dong-tai-tro-cac-du-an-cong-trinh-xanh}{bidv.com.vn (Green Credit News, Mar 2024)}
    \end{itemize}
\end{itemize}

\end{tcolorbox}


% =================== Box 5 — E (VPBank green credit 2024) ===================
\begin{tcolorbox}[colback=white,colframe=black!65,boxrule=0.6pt,arc=0pt,left=8pt,right=8pt,top=8pt,bottom=8pt]

\textbf{Question:} What was the total outstanding balance of green credit at VPBank as of December 31, 2024 (in VND billion)?

\textbf{Generated answer:} 21,943 billion VND. As of December 31, 2024, the total outstanding balance of green credit at VPBank was reported to be 21,943 billion VND, which represents a 48.5\% increase year-over-year. This information is sourced from VPBank's Sustainability Report 2024, page 23, and aligns with GRI standards GRI 3-3.

\begin{itemize}
    \item \textbf{Claim:} As of December 31, 2024, the total outstanding balance of green credit at VPBank was 21,943 billion VND.
    \begin{itemize}
        \item \textbf{Label:} SUPPORTS
        \item \textbf{Confidence:} 1.00
        \item \textbf{Source:} \href{https://static2.vietstock.vn/vietstock/2025/7/3/20250702___vpb___sustainability_report_2024.pdf}{vietstock.vn (VPBank Sustainability Report 2024)}
    \end{itemize}

    \item \textbf{Claim:} The green credit outstanding balance at VPBank increased by 48.5\% year-over-year in 2024.
    \begin{itemize}
        \item \textbf{Label:} SUPPORTS
        \item \textbf{Confidence:} 0.95
        \item \textbf{Source:} \href{https://amata.com/news-media-vietnam/93-articles-vietnam/1357-memorandum-of-understanding-between-amata-city-ha-long-and-vpbank}{amata.com (VPBank–Amata MoU)}
    \end{itemize}
\end{itemize}

\end{tcolorbox}

\end{document}
